\documentclass[11pt]{article}
\usepackage[margin=2cm]{geometry}

% Core packages
\usepackage{amsmath,amssymb}
\usepackage{tikz-cd}
\usepackage{multicol}
\usepackage[backend=biber,style=authoryear]{biblatex}
\addbibresource{bibliography.bib}

% Paragraphs
\setlength{\parindent}{0pt}
\setlength{\parskip}{1\baselineskip}

\title{Taller 2\\[0.5em]
\large Astrofísica Moderna\\
Universidad de Antioquia}
\author{Daniel Soto Villada}
\date{\today}

\begin{document}
\maketitle

\begin{enumerate}
    \item La radiofuente Cas A tiene una densidad de flujo de $3110\,\mathrm{Jy}$ a una frecuencia de $1\,\mathrm{GHz}$ en un intervalo de frecuencia (ancho de banda) de $100\,\mathrm{MHz}$.\\
    a) Si observamos esta fuente con un radiotelescopio de área colectora de $100\,\mathrm{m}^2$, ¿cuánta energía detectaremos en 1 segundo?\\
    b) Cas A es aproximadamente una fuente circular con un diámetro angular de unos 3 minutos de arco. ¿Cuál es la intensidad media de la fuente a una frecuencia de $1\,\mathrm{GHz}$?\\
    Nota: Dé sus respuestas en cgs.

    \textit{Solución.} Suponiendo que la densidad de flujo es aproximadamente constante en el ancho de banda, convertimos
    \[
        F_\nu=3110\,\mathrm{Jy}=3.11\times10^{-20}\,\mathrm{erg\,s^{-1}\,cm^{-2}\,Hz^{-1}}.
    \]
    Para $A=100\,\mathrm{m^2}=10^6\,\mathrm{cm^2}$, $\Delta\nu=100\,\mathrm{MHz}=10^8\,\mathrm{Hz}$ y $t=1\,\mathrm{s}$, la energía recibida es
    \[
        E=F_\nu\Delta\nu At=(3.11\times10^{-20})(10^8)(10^6)(1)
        =3.11\times10^{-6}\,\mathrm{erg}.
    \]
    El radio angular de Cas A es $1.5\,\mathrm{arcmin}=4.36\times10^{-4}\,\mathrm{rad}$, por lo que, aproximando la fuente como un disco uniforme, $\Omega\simeq\pi(4.36\times10^{-4})^2=5.98\times10^{-7}\,\mathrm{sr}$. Así, la intensidad media es
    \[
        \bar I_\nu=\frac{F_\nu}{\Omega}=5.20\times10^{-14}\,\mathrm{erg\,s^{-1}\,cm^{-2}\,Hz^{-1}\,sr^{-1}}.
    \]
    \parencite[Secs. 7.3.1.2, pp. 236--238; Sec. 8.1, pp. 255--260]{chaparro2026}

    \item Una estrella tiene magnitud aparente observada $m_{\mathrm{obs}} = 12.0$ y se encuentra a $500\,\mathrm{pc}$. Si el enrojecimiento interestelar provoca una extinción de $A_V = 0.8\,\mathrm{mag}$:
    \begin{itemize}
        \item Calcule la magnitud real (sin extinción):
        \[
            m_{\mathrm{real}} = m_{\mathrm{obs}} - A_V
        \]
        \item Halle su magnitud absoluta.
    \end{itemize}
    \textit{Solución.} Corregimos la magnitud aparente por extinción:
    \[
        m_{V,0}=m_{V,\mathrm{obs}}-A_V=12.0-0.8=11.2.
    \]
    Con el módulo de distancia, $m_{V,0}-M_V=5\log_{10}(d/10\,\mathrm{pc})$,
    \[
        M_V=11.2-5\log_{10}(500/10)=2.71\,\mathrm{mag}.
    \]
    \parencite[Sec. 7.3.3, pp. 241--248; Sec. 12.2.1, pp. 378--379]{chaparro2026}

    \item Una estrella tiene magnitudes aparentes $B = 8.75$ y $V = 7.90$, y se encuentra a una distancia de $100\,\mathrm{pc}$.
    \begin{itemize}
        \item Calcule el índice de color $(B - V)$.
        \item Estime la temperatura efectiva usando la aproximación que vimos en clase.
        \item Determine su magnitud absoluta $M_V$ y discuta su tipo espectral aproximado.
    \end{itemize}
    \textit{Solución.} El índice de color es
    \[
        B-V=8.75-7.90=0.85.
    \]
    Usando la aproximación color--temperatura del texto,
    \[
        T_{\mathrm{eff}}\simeq\frac{7090\,\mathrm{K}}{(B-V)+0.71}
        =\frac{7090}{1.56}\,\mathrm{K}\simeq4.54\times10^3\,\mathrm{K}.
    \]
    Suponiendo despreciable la extinción, el módulo de distancia para $100\,\mathrm{pc}$ es $5\log_{10}(100/10)=5$, de modo que
    \[
        M_V=7.90-5=2.90\,\mathrm{mag}.
    \]
    Esta temperatura corresponde aproximadamente a una estrella de tipo K tardío, alrededor de K4--K5. \parencite[Sec. 7.3.3, pp. 241--248; Sec. 11.1.2, pp. 337--340; Secs. 11.2 y 11.5.1, pp. 340--350, 360--361]{chaparro2026}

    \item Un rayo de radiación atraviesa un medio puramente absorbente con coeficiente de absorción constante $\kappa_\nu = 0.2\,\mathrm{cm}^{-1}$ y sin emisión ($j_\nu = 0$). La intensidad incidente en $s = 0$ es $I_\nu(0) = 5.0\,\mathrm{W\,m^{-2}\,Hz^{-1}\,sr^{-1}}$.
    \begin{itemize}
        \item Escriba la ecuación de transferencia radiativa:
        \[
            \frac{dI_\nu}{ds} = -\kappa_\nu I_\nu.
        \]
        \item Resuélvala para $I_\nu(s)$.
        \item Calcule $I_\nu$ después de atravesar 10 cm de material.
    \end{itemize}
    \textit{Solución.} Integramos la ecuación de transferencia con $j_\nu=0$ y la condición inicial dada:
    \[
        I_\nu(s)=I_\nu(0)e^{-\kappa_\nu s}.
    \]
    A los $10\,\mathrm{cm}$, $\tau_\nu=\kappa_\nu s=(0.2)(10)=2$, así que
    \[
        I_\nu(10\,\mathrm{cm})=5.0e^{-2}\simeq0.677\,\mathrm{W\,m^{-2}\,Hz^{-1}\,sr^{-1}}.
    \]
    \parencite[Secs. 13.3 y 13.7, pp. 391--405]{chaparro2026}

    \item Una estrella tiene una temperatura $T = 7000\,\mathrm{K}$. La observamos a través del medio interestelar (ISM), el cual produce una extinción de $A_B = 5\,\mathrm{mag}$ en la banda $B$ ($\lambda_B = 440\,\mathrm{nm}$). Supongamos que la ley de extinción sigue la dependencia
    \[
        A_\lambda \propto \lambda^{-1}.
    \]
    \begin{itemize}
        \item ¿Cuál sería la frecuencia (o longitud de onda) del pico observado del espectro, antes y después de considerar la extinción?
        \item ¿Cómo cambiaría el color aparente de la estrella respecto a su color intrínseco?
        \item ¿Cuál sería su magnitud de extinción en $V$ ($A_V$) según la ley dada?
        \item ¿Cuál sería el exceso de color $E(B - V)$?
        \item ¿A qué distancia aproximada podría encontrarse la estrella?
    \end{itemize}
    \textit{Solución.} Suponemos que el espectro intrínseco es aproximadamente de cuerpo negro. El máximo de la distribución por unidad de longitud de onda se obtiene con la ley de Wien:
    \[
        \lambda_{\max,\mathrm{int}}=\frac{2.898\times10^{-3}\,\mathrm{m\,K}}{7000\,\mathrm{K}}
        =414\,\mathrm{nm}.
    \]
    Después de la extinción, el flujo espectral observado es
    \[
        F_{\lambda,\mathrm{obs}}=B_\lambda(T)10^{-0.4A_\lambda},
        \qquad A_\lambda=5\frac{440\,\mathrm{nm}}{\lambda}.
    \]
    Definimos $x=hc/(\lambda k_BT)$; el máximo de esta distribución satisface
    \[
        \frac{5}{x}-\frac{1}{1-e^{-x}}-q=0,
        \qquad q=0.4\ln(10)A_B\frac{\lambda_B k_BT}{hc}\simeq0.986.
    \]
    La solución numérica es $x\simeq2.397$, que corresponde a $\lambda_{\max,\mathrm{obs}}\simeq857\,\mathrm{nm}$, en el infrarrojo cercano. La frecuencia $c/\lambda$ en el máximo de $F_\lambda$ cambia de $7.24\times10^{14}\,\mathrm{Hz}$ a $3.50\times10^{14}\,\mathrm{Hz}$. El máximo de $F_\nu$ no coincide con $c/\lambda_{\max}$: antes de la extinción está en $4.12\times10^{14}\,\mathrm{Hz}$ (raíz $x\simeq2.821$), mientras que después satisface $3/x-1/(1-e^{-x})-q=0$, con $x\simeq1.259$, y queda en $1.84\times10^{14}\,\mathrm{Hz}$.

    Tomando $\lambda_V\simeq550\,\mathrm{nm}$,
    \[
        A_V=A_B\frac{\lambda_B}{\lambda_V}=5\frac{440}{550}=4.0\,\mathrm{mag},
        \qquad E(B-V)=A_B-A_V=1.0\,\mathrm{mag}.
    \]
    Por tanto, $(B-V)_{\mathrm{obs}}=(B-V)_0+1.0\,\mathrm{mag}$: el color aparente es más rojo. Con el valor medio aproximado $a_V\simeq1\,\mathrm{mag\,kpc^{-1}}$, $d\simeq A_V/a_V\simeq4\,\mathrm{kpc}$. Esta distancia es solo orientativa, pues la cantidad de polvo varía con la línea de visión. \parencite[Secs. 9.4, pp. 286--287; Sec. 10.1, pp. 308--320; Secs. 12.2--12.3.1, pp. 375--385]{chaparro2026}
    \end{enumerate}

\printbibliography[title={Referencias}]

\end{document}
